We run four controlled variations. All use the same training draws as the main group (seeds 0--4) and the same test sets.

\subsection{Target transform (log target)}
The target has a heavy zero-inflated tail, and a monotone output nonlinearity is the main non-additive structure that global-epistasis models posit~\cite{otwinowski2018inferring}. To test whether the CNN's advantage comes from such a nonlinearity we retrain every system on $\log(y+0.01)$ (group \texttt{abl\_logtarget}; for the CNN, group \texttt{abl\_cnn}), with everything else unchanged; metrics are still computed on the raw fitness, so Spearman is comparable with the main table. Table~\ref{tab:logt} and Fig.~\ref{fig:logt} show the result.

\begin{table*}[t]\footnotesize\centering\renewcommand{\arraystretch}{0.8}
\caption{Log-target ablation (rand and two dbl sizes): Spearman, mean $\pm$ std over draws.}\label{tab:logt}
\resizebox{0.4\textwidth}{!}{\begin{tabular}{llcc}
\toprule
Method & Task & spearman $\uparrow$ & $n$ \\
\midrule
One-hot ridge log target & rand\_48 & 0.26 $\pm$ 0.05 & 5 \\
GP log target & rand\_48 & \underline{0.27 $\pm$ 0.03} & 5 \\
One-hot + pairwise ridge log target & rand\_48 & \textbf{0.27 $\pm$ 0.03} & 5 \\
\midrule
One-hot ridge log target & rand\_96 & 0.37 $\pm$ 0.05 & 5 \\
GP log target & rand\_96 & \underline{0.37 $\pm$ 0.04} & 5 \\
One-hot + pairwise ridge log target & rand\_96 & \textbf{0.38 $\pm$ 0.04} & 5 \\
\midrule
One-hot ridge log target & rand\_384 & \underline{0.52 $\pm$ 0.03} & 5 \\
GP log target & rand\_384 & \textbf{0.53 $\pm$ 0.03} & 5 \\
One-hot + pairwise ridge log target & rand\_384 & 0.52 $\pm$ 0.03 & 5 \\
\midrule
One-hot ridge log target & rand\_2000 & \underline{0.59 $\pm$ 0.01} & 5 \\
GP log target & rand\_2000 & \textbf{0.59 $\pm$ 0.01} & 5 \\
One-hot + pairwise ridge log target & rand\_2000 & 0.54 $\pm$ 0.01 & 5 \\
\midrule
One-hot ridge log target & dbl\_384 & \textbf{0.46 $\pm$ 0.02} & 5 \\
GP log target & dbl\_384 & \underline{0.46 $\pm$ 0.03} & 5 \\
One-hot + pairwise ridge log target & dbl\_384 & 0.44 $\pm$ 0.03 & 5 \\
\midrule
One-hot ridge log target & dbl\_2000 & \textbf{0.50 $\pm$ 0.00} & 5 \\
GP log target & dbl\_2000 & \underline{0.48 $\pm$ 0.01} & 5 \\
One-hot + pairwise ridge log target & dbl\_2000 & 0.45 $\pm$ 0.01 & 5 \\
\bottomrule
\end{tabular}
}
\end{table*}

\begin{figure}[t]\centering
\includegraphics[width=0.5\columnwidth]{logtarget_curves.pdf}
\caption{Mean Spearman versus $N$ (random training) with the raw target (solid) and the log-transformed target (dashed).}\label{fig:logt}
\end{figure}

With a log target every model improves substantially at $N\ge96$ in Spearman. At $N{=}384$ the four log-target systems lie within a narrow band (ridge \rhval{abl_logtarget/one-hot-ridge-log-target/rand_384/spearman/mean:3}, GP \rhval{abl_logtarget/gp-log-target/rand_384/spearman/mean:3}, pairwise \rhval{abl_logtarget/one-hot-+-pairwise-ridge-log-target/rand_384/spearman/mean:3}, CNN \rhval{abl_cnn/cnn-log-target/rand_384/spearman/mean:3}), whereas with the raw target the CNN led by a visible margin (\rhval{main/cnn-reimplemented/rand_384/spearman/mean:3} versus \rhval{main/one-hot-ridge-reimplemented/rand_384/spearman/mean:3}). At $N{=}2000$ the CNN remains slightly ahead (\rhval{abl_cnn/cnn-log-target/rand_2000/spearman/mean:3} versus \rhval{abl_logtarget/one-hot-ridge-log-target/rand_2000/spearman/mean:3} for ridge) and pairwise ridge falls behind (\rhval{abl_logtarget/one-hot-+-pairwise-ridge-log-target/rand_2000/spearman/mean:3}). The gap to the CNN closes at $N\le384$ and shrinks at $N{=}2000$ but does not vanish. This is consistent with part of the CNN's raw-target advantage being a learned output nonlinearity that a fixed transform provides to the linear models; it does not isolate epistasis from other effects, the CNN itself also gains from the log target, the log-target ablation was not run for the CNN on the dbl tasks, and the CNN could still exploit interactions at $N{=}2000$.
The transform has a cost: top-100 recall of additive ridge at $N{=}2000$ drops to \rhval{abl_logtarget/one-hot-ridge-log-target/rand_2000/top100_recall/mean:3} (raw target: \rhval{main/one-hot-ridge-reimplemented/rand_2000/top100_recall/mean:3}), because the logarithm compresses differences among high-fitness variants, so the better global ranking does not carry over to the tail. The other models keep most of their recall (CNN \rhval{abl_cnn/cnn-log-target/rand_2000/top100_recall/mean:3}, GP \rhval{abl_logtarget/gp-log-target/rand_2000/top100_recall/mean:3}).

\subsection{GP hyperparameters}
Replacing the marginal-likelihood grid by one fixed setting ($\sigma^2{=}1$, $\gamma{=}0.5$, $\sigma_n^2{=}0.1$; group \texttt{abl\_gp}) matters little at rand$_{384}$ but not at the larger sizes: rand$_{384}$ \rhval{abl_gp/gp-fixed-hyperparameters/rand_384/spearman/mean:3} (fitted: \rhval{main/gp-hamming-kernel-reimplemented/rand_384/spearman/mean:3}), rand$_{2000}$ \rhval{abl_gp/gp-fixed-hyperparameters/rand_2000/spearman/mean:3} (fitted \rhval{main/gp-hamming-kernel-reimplemented/rand_2000/spearman/mean:3}), dbl$_{2000}$ \rhval{abl_gp/gp-fixed-hyperparameters/dbl_2000/spearman/mean:3} (fitted \rhval{main/gp-hamming-kernel-reimplemented/dbl_2000/spearman/mean:3}). The effect is negligible at rand$_{384}$; at rand$_{2000}$ fitting raises the mean Spearman, whereas at dbl$_{2000}$ the fitted GP is lower than the fixed one, which is about as good as additive ridge (\rhval{main/one-hot-ridge-reimplemented/dbl_2000/spearman/mean:3}). Marginal-likelihood selection on low-order mutants may therefore pick hyperparameters that extrapolate poorly to higher-order variants; we did not run a test that isolates this mechanism, and the fixed setting was not tuned for any task.

\subsection{Ridge penalty}
Table~\ref{tab:alpha} fixes the penalty at each decade at rand$_{384}$ and compares with the CV choice. Additive ridge is almost flat in $\alpha$ (from \rhval{derived/derived-analysis/rand_384/alpha_ridge_0.001_mean/mean:3} at $10^{-3}$ to \rhval{derived/derived-analysis/rand_384/alpha_ridge_1000_mean/mean:3} at $10^{3}$), so CV matters little; pairwise ridge needs a large penalty (\rhval{derived/derived-analysis/rand_384/alpha_pairwise_0.001_mean/mean:3} at $10^{-3}$ versus \rhval{derived/derived-analysis/rand_384/alpha_pairwise_10_mean/mean:3} at $10$), and CV finds it (\rhval{derived/derived-analysis/rand_384/alpha_pairwise_cv_mean/mean:3}). The pairwise model's lack of benefit in Table~\ref{tab:main} is therefore not an artifact of a badly chosen penalty at this size.

\begin{table}[t]\footnotesize\centering\renewcommand{\arraystretch}{0.8}
\caption{Mean Spearman at rand$_{384}$ as a function of the fixed ridge penalty $\alpha$, and with the CV-selected penalty (main group).}\label{tab:alpha}
\resizebox{\columnwidth}{!}{\begin{tabular}{lrrrrrrrr}
\toprule
System & $10^{-3}$ & $10^{-2}$ & $10^{-1}$ & $1$ & $10$ & $10^2$ & $10^3$ & CV \\
\midrule
One-hot ridge & 0.392 & 0.392 & 0.392 & 0.393 & 0.394 & 0.385 & 0.380 & 0.381 \\
Pairwise ridge & 0.348 & 0.348 & 0.351 & 0.371 & 0.399 & 0.389 & 0.382 & 0.396 \\
\bottomrule
\end{tabular}
}
\end{table}

\subsection{CNN ensembling}
Averaging five CNNs (group \texttt{abl\_cnn}) gives rand$_{2000}$ Spearman \rhval{abl_cnn/cnn-ensemble-of-5/rand_2000/spearman/mean:3} against \rhval{main/cnn-reimplemented/rand_2000/spearman/mean:3} for the single network, and \rhval{abl_cnn/cnn-ensemble-of-5/rand_96/spearman/mean:3} against \rhval{main/cnn-reimplemented/rand_96/spearman/mean:3} at rand$_{96}$: ensembling helps the CNN further, so the single-network numbers in the main table understate what this architecture can do. Ensembles were not given to the other systems (an asymmetry we flag rather than resolve).
